\begin{theorem}[Linearity of the Lebesgue integral] \label{theorem:linearity_of_lebesgue_integral}
    \TODO{AI generated, verify}
    Let $(X,\mathcal{A},\mu)$ be a \CrefAndHyperrefIfExist{definition:sigma_algebra_on_a_set_and_measurable_space}{measure space}. If $f,g : X \to \mathbb{R}$ are \CrefAndHyperrefIfExist{definition:lebesgue_integral_of_measurable_functions_into_the}{(Lebesgue) integrable} with respect to $\mu$ and $\alpha,\beta \in \mathbb{R}$, then the function $\alpha f + \beta g : X \to \mathbb{R}$ is also Lebesgue integrable with respect to $\mu$, and
    $$\int_X (\alpha f(x) + \beta g(x))\,d\mu(x) = \alpha \int_X f(x)\,d\mu(x) + \beta \int_X g(x)\,d\mu(x).$$
\end{theorem}

\begin{proof}
    \TODO{AI generated, verify}
We use the construction of the Lebesgue integral in \CrefAndHyperrefIfExist{definition:lebesgue_integral_of_measurable_functions_into_the}: for a nonnegative measurable function $h$, the integral is defined as $$\int_X h\,d\mu = \sup\left\{\int_X s\,d\mu : 0 \leq s \leq h,\; s \text{ simple}\right\},$$ and for an integrable function $f: X \to \mathbb{R}$, the integral is $\int_X f\,d\mu = \int_X f^+\,d\mu - \int_X f^-\,d\mu$, where $f^+ = \max(f,0)$ and $f^- = \max(-f,0)$.

We proceed in three steps: first for nonnegative simple functions (direct computation), then for nonnegative measurable functions (supremum argument), and finally for general integrable functions (decomposition into positive and negative parts).

\begin{enumerate}
    \item We begin with nonnegative simple functions. Let $s,t : X \to [0,\infty)$ be simple functions, so that $$s = \sum_{i=1}^m a_i \mathbf{1}_{A_i}, \qquad t = \sum_{j=1}^n b_j \mathbf{1}_{B_j},$$ where $a_i,b_j \in [0,\infty)$ and $\{A_1,\dots,A_m\}$, $\{B_1,\dots,B_n\}$ are measurable partitions of $X$ (pairwise disjoint sets in $\mathcal{A}$ whose union is $X$). Let $\alpha,\beta \geq 0$.

    Consider the common refinement partition $\{C_{ij}\}_{i,j}$ of $X$, where $C_{ij} = A_i \cap B_j$. Each set $C_{ij}$ is measurable since it is an intersection of two measurable sets, and these sets are pairwise disjoint with union equal to $X$. On each $C_{ij}$, both $s$ and $t$ are constant: $$s(x) = a_i \quad\text{and}\quad t(x) = b_j \qquad\text{for all } x \in C_{ij}.$$ Therefore the function $\alpha s + \beta t$ is also simple, with representation
    $$\alpha s + \beta t = \sum_{i=1}^m \sum_{j=1}^n (\alpha a_i + \beta b_j) \mathbf{1}_{C_{ij}}.$$

    By the definition of the Lebesgue integral of a simple function in \CrefAndHyperrefIfExist{definition:lebesgue_integral_of_measurable_functions_into_the}, we compute
    $$\int_X (\alpha s + \beta t)\,d\mu = \sum_{i=1}^m \sum_{j=1}^n (\alpha a_i + \beta b_j) \mu(C_{ij}).$$

    We split this double sum:
    $$\sum_{i=1}^m \sum_{j=1}^n (\alpha a_i + \beta b_j)\mu(C_{ij}) = \alpha \sum_{i=1}^m \sum_{j=1}^n a_i \mu(C_{ij}) + \beta \sum_{i=1}^m \sum_{j=1}^n b_j \mu(C_{ij}).$$

    For the first term, we use that $\{C_{ij}\}_{j}$ is a partition of $A_i$ (since $\{B_j\}_j$ partitions $X$ and $A_i = \bigcup_j C_{ij}$), so $$\sum_{j=1}^n \mu(C_{ij}) = \mu(A_i).$$ Hence
    $$\alpha \sum_{i=1}^m a_i \left(\sum_{j=1}^n \mu(C_{ij})\right) = \alpha \sum_{i=1}^m a_i \mu(A_i) = \alpha \int_X s\,d\mu.$$
    Similarly for the second term: $\{C_{ij}\}_i$ partitions $B_j$, so $$\beta \sum_{j=1}^n b_j \left(\sum_{i=1}^m \mu(C_{ij})\right) = \beta \sum_{j=1}^n b_j \mu(B_j) = \beta \int_X t\,d\mu.$$

    This proves linearity for nonnegative simple functions with nonnegative coefficients:
    $$\int_X (\alpha s + \beta t)\,d\mu = \alpha \int_X s\,d\mu + \beta \int_X t\,d\mu.$$

    \item Now we extend to nonnegative measurable functions. Let $f,g : X \to [0,\infty]$ be measurable (not necessarily integrable; their integrals may be infinite) and let $\alpha,\beta \geq 0$. We must show that
    $$\int_X (\alpha f + \beta g)\,d\mu = \alpha \int_X f\,d\mu + \beta \int_X g\,d\mu.$$

    Since $f$ and $g$ are nonnegative measurable functions, there exist increasing sequences of nonnegative simple functions $(s_n)_n$ and $(t_n)_n$ such that $s_n \nearrow f$ and $t_n \nearrow g$ pointwise (by the standard approximation theorem for measurable functions). Then $(\alpha s_n + \beta t_n)_n$ is an increasing sequence of nonnegative simple functions converging pointwise to $\alpha f + \beta g$.

    By Step 1, each term satisfies $$\int_X (\alpha s_n + \beta t_n)\,d\mu = \alpha \int_X s_n\,d\mu + \beta \int_X t_n\,d\mu.$$

    We now apply the \CrefAndHyperrefIfExist{theorem:monotone_convergence_theorem_for_the_limit_of_lebesgue_integrals_of_an_increasing_sesquence_of_functions_and_the_lebesgue_integral_of_the_limit_of_the_functions}{Monotone Convergence Theorem} to pass limits. Since $s_n \nearrow f$, the MCT gives $\lim_{n\to\infty}\int_X s_n\,d\mu = \int_X f\,d\mu$. Similarly, $\lim_{n\to\infty}\int_X t_n\,d\mu = \int_X g\,d\mu$, and since $\alpha s_n + \beta t_n \nearrow \alpha f + \beta g$, the MCT gives $$\lim_{n\to\infty} \int_X (\alpha s_n + \beta t_n)\,d\mu = \int_X (\alpha f + \beta g)\,d\mu.$$

    Taking limits in the identity from Step 1:
    $$\int_X (\alpha f + \beta g)\,d\mu = \lim_{n\to\infty} \int_X (\alpha s_n + \beta t_n)\,d\mu = \lim_{n\to\infty}\left(\alpha \int_X s_n\,d\mu + \beta \int_X t_n\,d\mu\right).$$

    Since $\alpha,\beta \geq 0$, this equals
    $$\alpha \cdot \lim_{n\to\infty} \int_X s_n\,d\mu + \beta \cdot \lim_{n\to\infty} \int_X t_n\,d\mu = \alpha \int_X f\,d\mu + \beta \int_X g\,d\mu,$$
    where the limit of a sum equals the sum of limits (with the convention that $+\infty + c = +\infty$ for any nonnegative real or infinite value). This establishes linearity for nonnegative measurable functions with nonnegative coefficients.

    \item We now handle the general case where $f,g : X \to \mathbb{R}$ are integrable and $\alpha,\beta \in \mathbb{R}$ are arbitrary scalars. First we verify that $\alpha f + \beta g$ is integrable, i.e., $\int_X |\alpha f + \beta g|\,d\mu < \infty$. By the triangle inequality for real numbers and properties of absolute values (which follow from basic algebra in $\mathbb{R}$),
    $$|\alpha(x) f(x) + \beta g(x)| \leq |\alpha| \cdot |f(x)| + |\beta| \cdot |g(x)|.$$

    The function $x \mapsto |\alpha| \cdot |f(x)|$ is the product of a nonnegative constant and the nonnegative measurable function $|f| = f^+ + f^-$, so by Step 2 (linearity for nonnegative functions),
    $$\int_X |\alpha|\cdot |f|\,d\mu = |\alpha| \int_X |f|\,d\mu < \infty,$$
    since $f$ is integrable. Similarly $\int_X |\beta|\cdot |g|\,d\mu < \infty$. By Step 2 again (additivity for nonnegative functions), the sum $|\alpha||f| + |\beta||g|$ has finite integral. Since $|\alpha f + \beta g|$ is dominated by this function, its integral is also finite. Hence $\alpha f + \beta g$ is integrable.

    To establish the integral identity, we use the decomposition into positive and negative parts. For any real-valued function $h$, write $h = h^+ - h^-$. We compute $(\alpha f + \beta g)^+$ and $(\alpha f + \beta g)^-$ in terms of $f^\pm$ and $g^\pm$.

    We consider two cases based on the signs of $\alpha$ and $\beta$. By symmetry, it suffices to handle the case where both are nonnegative; the general case follows by writing negative scalars as negatives of positive ones.

    First assume $\alpha,\beta \geq 0$. Then we claim that $$(\alpha f + \beta g)^+ = (\alpha f^+ + \beta g^+) - (\alpha f^- + \beta g^-),$$ and moreover the two terms on the right are nonnegative. To see this, note that $\alpha f + \beta g = \alpha(f^+ - f^-) + \beta(g^+ - g^-) = (\alpha f^+ + \beta g^+) - (\alpha f^- + \beta g^-)$ by ordinary algebra in $\mathbb{R}$. Let $P = \alpha f^+ + \beta g^+$ and $N = \alpha f^- + \beta g^-$. Both are nonnegative measurable functions. The function $P - N$ has $(P-N)^+ = P$ and $(P-N)^- = N$ only when $PN = 0$ everywhere, which is not generally true. So we cannot simply separate them this way.

    Instead, we use a more direct approach. For any real numbers $u,v$ and $\alpha,\beta \geq 0$, the identity $(\alpha u + \beta v)^+ - (\alpha u + \beta v)^- = \alpha(u^+ - u^-) + \beta(v^+ - v^-)$ holds as an equality in $\mathbb{R}$. Integrating both sides and using linearity for nonnegative functions:
    $$\int_X (\alpha f + \beta g)\,d\mu := \int_X (\alpha f + \beta g)^+\,d\mu - \int_X (\alpha f + \beta g)^-\,d\mu.$$

    We relate $(\alpha f + \beta g)^\pm$ to $f^\pm$ and $g^\pm$. For arbitrary real numbers $u,v$, we have the identity
    $$\alpha(u^+ - u^-) + \beta(v^+ - v^-) = (\alpha u + \beta v)^+ - (\alpha u + \beta v)^-.$$

    We decompose each of $\alpha f$ and $\beta g$ using the fact that $(\alpha f)^+ = |\alpha|f^+$ and $(\alpha f)^- = |\alpha|f^-$. Indeed, if $u \geq 0$, then $(\alpha u)^+ = \alpha u = |\alpha|u = |\alpha|\cdot u^+$ and $(\alpha u)^- = 0 = |\alpha|\cdot u^-$. If $u < 0$, then $\alpha u \leq 0$ (for $\alpha \geq 0$), so $(\alpha u)^+ = 0 = |\alpha|\cdot u^+$ and $(\alpha u)^- = -\alpha u = \alpha(-u) = |\alpha|u^-$. In both cases, the identity holds.

    Therefore we can write $$\alpha f + \beta g = (\alpha f^+ - \alpha f^-) + (\beta g^+ - \beta g^-).$$ Let $A = \alpha f^+ + \beta g^+$ and $B = \alpha f^- + \beta g^-$. Both are nonnegative measurable functions, so by Step 2 they have well-defined integrals. We have $\alpha f + \beta g = A - B$, but this is not the same as $(\alpha f+\beta g)^+ - (\alpha f+\beta g)^-$ unless $AB = 0$.

    To resolve this, we use the following identity valid for any real numbers $x,y$:
    $$(x+y)^+ + y = (x-y)^+ + x.$$ This is a standard algebraic identity. Setting $x = \alpha f - \beta g$ and using the decomposition more carefully:

    Actually, let us use an even simpler route. By definition, $\int_X (\alpha f + \beta g)\,d\mu = \int_X (\alpha f + \beta g)^+\,d\mu - \int_X (\alpha f + \beta g)^-\,d\mu$. We prove that
    $$\int_X (\alpha f + \beta g)^+\,d\mu = \alpha \int_X f^+\,d\mu + \beta \int_X g^+\,d\mu$$ and similarly for the negative parts. This would follow if $(\alpha f + \beta g)^+ = \alpha f^+ + \beta g^+$, but this is false in general (e.g., $f=1,g=-1,\alpha=\beta=1$ gives $(f+g)^+=0$ while $f^++g^+=2$).

    Therefore we use the approach of decomposing into positive and negative parts more carefully. We write $\alpha f + \beta g = (\alpha f^+ + \beta g^+) - (\alpha f^- + \beta g^-)$. Let $P = \alpha f^+ + \beta g^+$ and $Q = \alpha f^- + \beta g^-$. Then $(\alpha f + \beta g)^+ - (\alpha f + \beta g)^- = P - Q$, which means $$(\alpha f + \beta g)^+ + Q = (\alpha f + \beta g)^- + P.$$

    All four functions here are nonnegative measurable. By Step 2 (linearity for nonnegative functions, including additivity), we can integrate both sides:
    $$\int_X (\alpha f + \beta g)^+\,d\mu + \int_X Q\,d\mu = \int_X (\alpha f + \beta g)^-\,d\mu + \int_X P\,d\mu.$$

    By Step 2 again, $\int_X P\,d\mu = \alpha \int_X f^+\,d\mu + \beta \int_X g^+\,d\mu$ and $\int_X Q\,d\mu = \alpha \int_X f^-\,d\mu + \beta \int_X g^-\,d\mu$. Substituting:
    $$\int_X (\alpha f + \beta g)^+\,d\mu + \alpha \int_X f^-\,d\mu + \beta \int_X g^-\,d\mu = \int_X (\alpha f + \beta g)^-\,d\mu + \alpha \int_X f^+\,d\mu + \beta \int_X g^+\,d\mu.$$

    Rearranging:
    $$\int_X (\alpha f + \beta g)^+\,d\mu - \int_X (\alpha f + \beta g)^-\,d\mu = \alpha\left(\int_X f^+\,d\mu - \int_X f^-\,d\mu\right) + \beta\left(\int_X g^+\,d\mu - \int_X g^-\,d\mu\right).$$

    The left side is $\int_X (\alpha f + \beta g)\,d\mu$ by definition. The right side is $\alpha \int_X f\,d\mu + \beta \int_X g\,d\mu$. All terms are finite since $f,g$ are integrable, so no indeterminate forms arise.

    For the case where $\alpha < 0$, write $\alpha = -|\alpha|$. Then $(\alpha f)^+ = |\alpha|f^-$ and $(\alpha f)^- = |\alpha|f^+$ (since $-\lambda u \geq 0$ iff $u \leq 0$ for $\lambda > 0$). By the same argument as above applied to $(-|\alpha|f) + (\beta g)$, we obtain
    $$\int_X (-|\alpha|f + \beta g)\,d\mu = -|\alpha| \int_X f\,d\mu + \beta \int_X g\,d\mu = \alpha \int_X f\,d\mu + \beta \int_X g\,d\mu.$$

    The case $\beta < 0$ is handled identically. For arbitrary signs of $\alpha,\beta$, we apply the result step by step: first for scalar multiplication (which gives $\int (\alpha f) = \alpha \int f$), then for addition (which gives $\int(f+g) = \int f + \int g$). Thus
    $$\int_X (\alpha f + \beta g)\,d\mu = \alpha \int_X f\,d\mu + \beta \int_X g\,d\mu.$$

    This completes the proof of linearity for the Lebesgue integral.
\end{enumerate}
This completes the proof of linearity for the Lebesgue integral.
\end{proof}
